\ours{} needs categorical factors whose levels an LLM can recognize by name or structure; on a purely continuous problem it has nothing to ask and reduces to a GP. The number of possible comparisons grows with the square of the number of reagents. Restricting candidates to the top experiments under the acquisition keeps the number asked small, but scoring them on pools with thousands of reagents per factor will need amortization. The query effect in Eq.~\eqref{eq:probit} absorbs errors shared by repeats of one question, while errors that correlate across different questions, such as a systematic misconception about a ligand class, still violate assumption (A3) of Proposition~\ref{prop:consistency}; the counterfactual and inverted oracles bound their cost empirically but do not model them. Axis discovery chooses what to ask from the same experiments it is later judged on; excluding the shown reagents from the new axis's comparisons removes the most direct reuse but not all of it. Finally, all campaigns are retrospective on pools of measured or emulated reactions, and a prospective campaign in a laboratory remains to be run.
